\documentclass[a4paper,10pt]{article}
\usepackage{lamia-pedagogy}

\begin{document}
% PAGE 4 : Compétence 03 - Compréhension de l'écrit (Phase 1)
% =========================================================================
\ficheHeader{1 Heure}{\makebox[2.5cm]{\dotfill}}{1\textsuperscript{ère} AM}{2026 / 2027}

\vspace{-2pt}
\competenceBanner{\faBookReader\ Compétence 03 : Compréhension de l'écrit}

\begin{minipage}[t]{0.48\textwidth}
\begin{tcolorbox}[
  enhanced,
  colback=secondarylight,
  colframe=secondary,
  arc=1.5mm,leftrule=3pt,rightrule=0.5pt,toprule=0.5pt,bottomrule=0.5pt,
  top=3pt,bottom=3pt,left=6pt,right=6pt,
  title={\small\textbf{\faBullseye\ Critères 1 et 2}},
  coltitle=secondary!90!black,attach boxed title to top left={yshift=-2mm,xshift=4mm},
  boxed title style={colback=white,boxrule=0.5pt,colframe=secondary,arc=1mm,top=1pt,bottom=1pt}
]
\vspace{1pt}
{\scriptsize
Identifier le thème général et le champ lexical relatif au texte.
\begin{itemize}[leftmargin=*,noitemsep,topsep=1pt]
  \item[\faCheckCircle] Repérer les indices textuels thématiques.
  \item[\faCheckCircle] Sélectionner le vocabulaire lié au domaine.
\end{itemize}
}
\end{tcolorbox}
\end{minipage}
\hfill
\begin{minipage}[t]{0.48\textwidth}
\begin{tcolorbox}[
  enhanced,
  colback=accentlight,
  colframe=accent,
  arc=1.5mm,leftrule=3pt,rightrule=0.5pt,toprule=0.5pt,bottomrule=0.5pt,
  top=3pt,bottom=3pt,left=6pt,right=6pt,
  title={\small\textbf{\faExclamationTriangle\ Difficultés identifiées}},
  coltitle=accent!90!black,attach boxed title to top left={yshift=-2mm,xshift=4mm},
  boxed title style={colback=white,boxrule=0.5pt,colframe=accent,arc=1mm,top=1pt,bottom=1pt}
]
\vspace{1pt}
{\scriptsize
\begin{itemize}[leftmargin=*,noitemsep,topsep=1pt]
  \item[\faTimesCircle] L'élève peine à repérer le thème central.
  \item[\faTimesCircle] Difficulté à regrouper les mots d'un même champ lexical.
\end{itemize}
}
\end{tcolorbox}
\end{minipage}

\vspace{1pt}
\begin{tcolorbox}[
  enhanced,
  colback=secondarylight,
  colframe=secondary,
  arc=1.5mm,leftrule=3.5pt,rightrule=0.6pt,toprule=0.6pt,bottomrule=0.6pt,
  top=2.5pt,bottom=2.5pt,left=8pt,right=8pt
]
{\small
\textbf{\faUsers\ Modalité de travail :} En groupe puis individuel.
}
\end{tcolorbox}

\vspace{1pt}
\begin{tcolorbox}[
  enhanced,
  colback=white,
  colframe=cardborder,
  arc=1.5mm,
  top=2pt,bottom=2pt,left=6pt,right=6pt,
  title={\scriptsize\textbf{\faUsers\ Liste des élèves concernés}},
  coltitle=primary,attach boxed title to top left={yshift=-1.8mm,xshift=3mm},
  boxed title style={colback=primarylight,boxrule=0.4pt,colframe=cardborder,arc=1mm,top=0.5pt,bottom=0.5pt}
]
\vspace{1pt}
{\scriptsize
\begin{tabularx}{\linewidth}{@{}l X l X@{}}
1. & \makebox[\linewidth]{\dotfill} & 6. & \makebox[\linewidth]{\dotfill} \\
2. & \makebox[\linewidth]{\dotfill} & 7. & \makebox[\linewidth]{\dotfill} \\
3. & \makebox[\linewidth]{\dotfill} & 8. & \makebox[\linewidth]{\dotfill} \\
4. & \makebox[\linewidth]{\dotfill} & 9. & \makebox[\linewidth]{\dotfill} \\
5. & \makebox[\linewidth]{\dotfill} & 10. & \makebox[\linewidth]{\dotfill} \\
\end{tabularx}
}
\end{tcolorbox}

\vspace{1pt}
\textbf{\faComments\ I) Mise en situation :}\\
{\small
$\bullet$ Proposer aux élèves de citer quelques noms de métiers.\\
$\bullet$ Que porte un médecin ? \textit{Une blouse blanche (tenue médicale) et du matériel médical.}
}

\vspace{2pt}
\textbf{\faEye\ II) Phase d'identification :}\\
{\small\textit{Proposer aux élèves les trois textes suivants sans titres, puis leur demander de les associer aux titres et images :}}

\vspace{2pt}
\begin{tcolorbox}[
  enhanced,
  colback=blue!3!white,
  colframe=primary!60,
  arc=1.5mm,leftrule=3.5pt,rightrule=0.5pt,toprule=0.5pt,bottomrule=0.5pt,
  top=3pt,bottom=3pt,left=7pt,right=7pt,
  title={\small\textbf{Texte 1}},
  coltitle=primary!90!black,attach boxed title to top left={yshift=-1.8mm,xshift=3mm},
  boxed title style={colback=white,boxrule=0.4pt,colframe=primary!60,arc=1mm,top=0.5pt,bottom=0.5pt}
]
\vspace{1pt}
{\small
Le pompier intervient lors des incendies pour éteindre le feu et sauver les personnes en danger. Il utilise des tuyaux, des lances à eau et porte une tenue spéciale pour se protéger. Le pompier aide aussi lors des accidents et des inondations. Le métier de pompier est très noble car il permet de sauver des vies.
}
\end{tcolorbox}

\vspace{-2pt}
\begin{tcolorbox}[
  enhanced,
  colback=emerald!5!white,
  colframe=secondary,
  arc=1.5mm,leftrule=3.5pt,rightrule=0.5pt,toprule=0.5pt,bottomrule=0.5pt,
  top=3pt,bottom=3pt,left=7pt,right=7pt,
  title={\small\textbf{Texte 2}},
  coltitle=secondary!90!black,attach boxed title to top left={yshift=-1.8mm,xshift=3mm},
  boxed title style={colback=white,boxrule=0.4pt,colframe=secondary,arc=1mm,top=0.5pt,bottom=0.5pt}
]
\vspace{1pt}
{\small
Le lapin est un animal domestique. Il est doux et apprécié des enfants. Il se nourrit de carottes, de légumes et d'herbe. Il vit dans une cage ou dans un clapier et peut être élevé à la ferme ou dans la maison.
}
\end{tcolorbox}

\vspace{-2pt}
\begin{tcolorbox}[
  enhanced,
  colback=amber!5!white,
  colframe=accent,
  arc=1.5mm,leftrule=3.5pt,rightrule=0.5pt,toprule=0.5pt,bottomrule=0.5pt,
  top=3pt,bottom=3pt,left=7pt,right=7pt,
  title={\small\textbf{Texte 3}},
  coltitle=accent!90!black,attach boxed title to top left={yshift=-1.8mm,xshift=3mm},
  boxed title style={colback=white,boxrule=0.4pt,colframe=accent,arc=1mm,top=0.5pt,bottom=0.5pt}
]
\vspace{1pt}
{\small
L'environnement est le lieu qui nous entoure. Pour vivre sainement, nous devons protéger ce milieu de tout ce qui peut le polluer. Pour cela, la participation de tous les citoyens est nécessaire.
}
\end{tcolorbox}

\vspace{1pt}
{\small
\textbf{\faLink\ Consigne : Relie chaque texte à son titre et au mot-clé correspondant :}\\
\vspace{1pt}
\begin{tabularx}{\linewidth}{XXX}
\badge[primary]{* Le pompier} & \badge[secondary]{* Le lapin} & \badge[accent]{* L'environnement} \\
\badge[primary]{* Sauver} & \badge[secondary]{* Domestique} & \badge[accent]{* Polluer} \\
\end{tabularx}
}

\vspace{3pt}
\begin{minipage}{0.32\textwidth}
\centering
\begin{tcolorbox}[enhanced,colback=white,colframe=cardborder,arc=2mm,top=2pt,bottom=2pt,left=2pt,right=2pt]
\centering
\includegraphics[height=3.4cm,keepaspectratio]{images/pompier.png}\\[2pt]
{\scriptsize\textbf{\faFireExtinguisher\ Le Pompier}}
\end{tcolorbox}
\end{minipage}
\hfill
\begin{minipage}{0.32\textwidth}
\centering
\begin{tcolorbox}[enhanced,colback=white,colframe=cardborder,arc=2mm,top=2pt,bottom=2pt,left=2pt,right=2pt]
\centering
\includegraphics[height=3.4cm,keepaspectratio]{images/lapin.png}\\[2pt]
{\scriptsize\textbf{\faPaw\ Le Lapin}}
\end{tcolorbox}
\end{minipage}
\hfill
\begin{minipage}{0.32\textwidth}
\centering
\begin{tcolorbox}[enhanced,colback=white,colframe=cardborder,arc=2mm,top=2pt,bottom=2pt,left=2pt,right=2pt]
\centering
\includegraphics[height=3.4cm,keepaspectratio]{images/environnement.png}\\[2pt]
{\scriptsize\textbf{\faLeaf\ L'Environnement}}
\end{tcolorbox}
\end{minipage}

\newpage
% =========================================================================
% PAGE 5 : Compétence 03 - Suite (Application & Réinvestissement)
% =========================================================================

\vspace*{2pt}
\badge[secondary]{\faTasks\ III) Phase d'application}

\vspace{3pt}
{\small\textbf{Activité 01 : Choisir le titre approprié}\\
\textit{Lis le texte suivant puis réponds aux questions :}}

\begin{tcolorbox}[
  enhanced,
  colback=amber!4!white,
  colframe=accent!80,
  arc=2mm,leftrule=3.5pt,rightrule=0.5pt,toprule=0.5pt,bottomrule=0.5pt,
  top=5pt,bottom=5pt,left=8pt,right=8pt,
  title={\small\textbf{\faBookOpen\ Texte :}},
  coltitle=accent!90!black,attach boxed title to top left={yshift=-2mm,xshift=4mm},
  boxed title style={colback=white,boxrule=0.4pt,colframe=accent!80,arc=1mm,top=1pt,bottom=1pt}
]
\vspace{1pt}
{\small
La girafe est un animal sauvage. Elle est le plus grand animal terrestre. Elle se nourrit de feuilles et de fruits qu'elle cueille en haut des arbres grâce à son long cou. Elle vit dans la savane et se déplace en troupeaux.
}
\end{tcolorbox}

\vspace{3pt}
{\small
\textbf{1/ Quel titre convient le mieux à ce texte ?}\par\vspace{2pt}
\noindent\begin{tabular}{@{\hspace{2mm}}p{4.5cm} p{4.5cm} p{6cm}@{}}
\qcmCircle \textbf{a)} Le lapin & \qcmCircle \textbf{b)} La girafe & \qcmCircle \textbf{c)} Le tremblement de terre \\
\end{tabular}\par\vspace{3pt}

\textbf{2/ Relève du texte le champ lexical relatif au mot \guillemotleft\ Animal\ \guillemotright\ :}\par\vspace{2pt}
\noindent\makebox[\linewidth]{\hspace*{2mm}\textbf{Réponse :} \dotfill}
}

\vspace{4pt}
\badge[primary]{\faPencilRuler\ IV) Phase de réinvestissement}

\vspace{2pt}
{\small
\textbf{Activité 02 : Complète le texte à l'aide de la boîte à mots :}\\
\centerline{\badge[accent]{\textbf{animaux}} \quad \badge[accent]{\textbf{grimaces}} \quad \badge[accent]{\textbf{girafes}} \quad \badge[accent]{\textbf{pain}}}

\vspace{3pt}
\begin{tcolorbox}[enhanced,colback=secondarylight,colframe=secondary,arc=2mm,top=4pt,bottom=4pt,left=8pt,right=8pt]
{\small
J'ai visité le Zoo avec ma famille. J'ai vu beaucoup d'\makebox[2.4cm]{\dotfill}\ . Les singes m'ont beaucoup amusé : ils faisaient des \makebox[2.4cm]{\dotfill}\ . J'ai donné à manger aux \makebox[2.4cm]{\dotfill}\ . \guillemotleft\ Il est interdit de donner du \makebox[2cm]{\dotfill}\ aux animaux !\ \guillemotright\ a rappelé le gardien.
}
\end{tcolorbox}
}

\vspace{4pt}
\begin{minipage}{0.48\textwidth}
\centering
\begin{tcolorbox}[enhanced,colback=white,colframe=cardborder,arc=2.5mm,top=3pt,bottom=3pt,left=3pt,right=3pt]
\centering
\includegraphics[width=\linewidth,height=8.8cm,keepaspectratio]{images/girafe_arbre.png}\\[3pt]
{\scriptsize\textbf{\faTree\ La girafe cueillant les feuilles en haut des arbres}}
\end{tcolorbox}
\end{minipage}
\hfill
\begin{minipage}{0.48\textwidth}
\centering
\begin{tcolorbox}[enhanced,colback=white,colframe=cardborder,arc=2.5mm,top=3pt,bottom=3pt,left=3pt,right=3pt]
\centering
\includegraphics[width=\linewidth,height=8.8cm,keepaspectratio]{images/girafe_savane.png}\\[3pt]
{\scriptsize\textbf{\faSun\ La girafe dans son habitat naturel (la savane)}}
\end{tcolorbox}
\end{minipage}

\newpage

\end{document}
